\documentclass[UTF8,11pt]{ctexart}
\usepackage[margin=1in]{geometry}
\usepackage{amsmath,amssymb,booktabs}
\usepackage[hidelinks]{hyperref}

\title{面向人员运动预测的可学习随机微分方程}
\author{作者信息待审批后补充}
\date{草稿，不包含实证结论}

\begin{document}
\maketitle

\begin{abstract}
搜救场景中的人员运动预测不能只给出单一路径，还应描述未来状态的不确定性。本文草稿围绕可学习随机微分方程模型、相应的 Python 实现和独立维护的 Lean 形式化工程组织论文主体。现有可执行基线是段级隐状态欠阻尼 Langevin 模型，支持精确高斯、Euler--Maruyama、分裂步进和公共随机数推断。Lean 工程验证了软证据与混合权重的一组代数恒等式。本文不报告私有轨迹数据、训练检查点、实验指标或性能结论。实验部分只给出待审阅的方案，后续仅可根据获准数据和版本化运行记录补充汇总证据。
\end{abstract}

\section{引言}

搜救规划需要在多个预测时域上获得人员可能位置和速度的分布。观测不完整、行为差异和环境条件都会影响预测，因此模型需要保留不确定性。随机微分方程可以描述状态的连续演化，并通过样本或矩给出预测分布。

本文如实记录现有研究资产及其边界。Python 仓库实现了模型拟合与推断框架，当前基线在每个轨迹段内采用固定的隐状态。Lean 仓库形式化了软证据的有限维、有限测度代数部分。抛物型偏微分方程的解存在性与唯一性，以及 Doob \(h\)-变换的构造，不在该形式化工程的范围内。

待补充实证材料后，本文将形成三项有限的工作内容：说明可审计的软件实现，划定机器检查定理的边界，并给出预先确定的评估方案。当前版本不对搜救效果或预测精度作出结论。

\section{问题表述}

设 \(X_t=(p_t,v_t)^\top\in\mathbb{R}^2\) 表示一个运动段在单个坐标方向上的位置和速度状态。带上下文的模型可写为

\begin{equation}
  dX_t=b_\theta(t,X_t,c_t)\,dt+\sigma_\theta(t,X_t,c_t)\,dW_t,
  \label{eq:general-sde-zh}
\end{equation}

其中，\(c_t\) 为经批准的上下文输入，\(W_t\) 为 Wiener 过程。当前公开的软件仓库没有注册可用的神经随机微分方程，也没有公开的上下文数据管线。因此，式\eqref{eq:general-sde-zh} 只给出建模接口，不表示已经完成的实证模型。

给定初始状态后，预测任务在若干时域 \(\tau\) 上返回 \(X_{t+\tau}\) 的样本或矩。评估应使用仓库外维护的划分文件，且只发布经过审查的汇总统计量。

\section{已实现的段级模型}

当前基线是带隐状态 \(r\in\{1,\ldots,K\}\) 的仿射欠阻尼 Langevin 随机微分方程。隐状态由模型上下文显式给出，并在单个运动段内保持不变。模型为

\begin{align}
  dp_t &= v_t\,dt, \\
  dv_t &= \left[-\Gamma_r v_t+(a_r-\kappa)p_t+c_r\right]dt+g_r\,dW_t.
  \label{eq:segment-sde-zh}
\end{align}

\(\Gamma_r\)、\(a_r\)、\(c_r\) 和 \(g_r\) 为状态相关参数，\(\kappa\) 在各状态之间共享。该模型的状态转移为仿射高斯形式，因此实现中提供了对应的精确转移核。

估计代码根据速度和位移变化构造特征，对各段做聚类初始化。随后交替计算段后验权重和加权最小二乘转移更新，再将离散转移转换为连续时间参数。这一描述对应软件实现，不能替代其在目标数据集上的收敛性验证。

\section{推断与软件边界}

Python 框架将随机微分方程模型、估计器、推断引擎、评价规则和结果适配器分离。能够提供精确转移核的模型可使用精确高斯采样。Euler--Maruyama 是显式选择的数值推断方式。段级模型还提供分裂步进方法，公共随机数可用于共享路径的多时域采样。不支持的模型与推断组合会直接报错，不会隐式切换为其他算法。

仓库中的合成数据流程是确定性的，用于安装与回归测试，不构成真实数据上的预测证据。神经随机微分方程和 Fokker--Planck 变体仍是未注册的研究骨架，条件特征对齐也尚未完全进入声明式数据源层。

\section{软证据与形式化边界}

对先验概率为 \(\pi\) 的二元软证据，记两类似然为 \(\ell_1\) 与 \(\ell_0\)，后验混合权重为

\begin{equation}
  w=\frac{\ell_1\pi}{\ell_1\pi+\ell_0(1-\pi)}.
  \label{eq:posterior-weight-zh}
\end{equation}

Lean 工程验证了该混合恒等式、式\eqref{eq:posterior-weight-zh} 的似然比形式、混合分布的总变差界和 logistic 权重的若干性质。它还验证了证据模式权重的乘积恒等式，以及在定理签名所列假设下的嵌套证据集中结论。

这些结论不构成目标随机微分方程适定性或 Doob \(h\)-变换的形式化证明。论文中只有 Lean 源码和公理审计实际覆盖的声明才能称为机器检查结果。

\section{数据治理与可复现性}

论文仓库不得包含原始或变换后的轨迹、坐标、时间戳、标识符、逐样本预测、非公开检查点，或含有上述内容的压缩文件。现有软件契约要求从受控的本地数据根读取主轨迹表和数据划分文件。该权威物理位置及其数据保管人尚未由本文草稿确认。

每项结果都应记录逻辑数据版本、划分文件版本、Python 与 Lean 仓库的提交版本、随机种子策略、模型配置和只含汇总信息的输出。受控运行器可对批准的数据根做只读挂载。公开持续集成只运行确定性的合成数据。

\section{实验方案}

本节是实验方案，不是结果报告。在添加任何表格或结论前，应固定并审查数据快照、划分文件、许可和隐私等级；预测时域、观测窗口、状态构造和缺失数据处理；比较模型及其推断配置；汇总评分规则和能量评分的具体约定；随机种子数量、区间估计和失败运行的处理规则；以及对已知数值问题和方案偏离的披露方式。

表\ref{tab:results-template-zh} 为待填模板。只有经批准的运行产生数值，并且数值与记录的配置相互核对后，才可填写其中的单元格。

\begin{table}[ht]
\centering
\caption{汇总评估结果模板，当前不填入数值。}
\label{tab:results-template-zh}
\begin{tabular}{llll}
\toprule
模型 & 时域 & 汇总评分 & 运行记录 \\
\midrule
段级随机微分方程 & [待确定] & [尚无结果] & [待确定] \\
批准的基线模型 & [待确定] & [尚无结果] & [待确定] \\
\bottomrule
\end{tabular}
\end{table}

\section{局限性与结论}

现有资产提供了可执行的基线模型、明确的数值推断路由、确定性的合成测试和范围受限的形式化证明。它们尚不足以说明私有轨迹上的预测表现、搜救任务中的业务价值，或 Lean 范围以外的解析性质。后续版本只能加入经过审查的汇总证据；软件或形式化边界发生变化时，也应同步修改本节。

\paragraph{源代码说明。} 本文草稿中的技术表述对应
\href{https://github.com/HuAzurestar/learnable-sde-for-movement-prediction}{Python 实现仓库}
和 \href{https://github.com/HuAzurestar/learnable-sde-theory-to-predict-movement}{Lean 形式化仓库}。
文献条目尚未核验，不能虚构；提交前应将核验后的条目写入 \texttt{references.bib}。

\end{document}
